%\begin{enumerate}[label=\arabic*.,ref=\thesection]
%\begin{enumerate}[label=\arabic*.,ref=\thesection.\theenumi]
%\begin{enumerate}[label=\thesection.\arabic*,ref=\thesection.\theenumi]
\renewcommand{\theequation}{\theenumi}

\begin{enumerate}[label=\arabic*.,ref=\thesubsection.\theenumi]

\item The {\em inner product} of  $\vec{P}$ and $\vec{Q}$ is defined as
%\numberwithin{equation}{enumi}
\begin{align}
\vec{P}^T\vec{Q} = p_1q_1+p_2q_2
\end{align}
\item The {\em norm} of a vector 
\begin{align}
\vec{P} = \myvec{p_1\\p_2}
\end{align}
is defined as
\begin{align}
\norm{\vec{P}} = \sqrt{p_1^2+p_2^2}
\end{align}
\item The {\em length} of $PQ$ is defined as
\begin{align}
\norm{\vec{P}-\vec{Q}}
\end{align}

%\renewcommand{\theequation}{\theenumi}

%
\item The {\em direction vector} of the line $PQ$ is defined as 
\begin{align}
\vec{P}-\vec{Q} = \myvec{p_1-q_1\\p_2-q_2}
\end{align}
%
%\item 
%\begin{equation}
%PQ \perp RS
%\iff \brak{\vec{P}-\vec{Q}}^T\brak{\vec{R}-\vec{S}} = 0
%\end{equation}
\item The point dividing   $PQ$  in the ratio $k:1$ is
\begin{equation}
\vec{R} = \frac{k\vec{P}+\vec{Q}}{k+1}
\end{equation}
%
\item The {\em area} of $\triangle PQR$ is the {\em determinant}
\begin{equation}
\begin{vmatrix}
1 & 1 & 1
\\
\vec{P} & \vec{Q} &\vec{R}
\end{vmatrix}
\end{equation}
\item {\em Orthogonality:} See Fig. \ref{fig:orth}.  In $\triangle ABC, AB \perp BC$. Show that
\begin{equation}
\brak{\vec{A}-\vec{B}}^T\brak{\vec{B}-\vec{C}} = 0
\label{eq:orth}
\end{equation}
\begin{figure}
\centering
\includegraphics[width=\columnwidth]{./line/figs/orth.eps}
\caption{}
\label{fig:orth}
\end{figure}
%
\solution Using Baudhayana's theorem,
\begin{align}
\norm{\vec{A}-\vec{B}}^2 + \norm{\vec{B}-\vec{C}}^2 &= 
\norm{\vec{C}-\vec{A}}^2
\\
\implies 
\brak{\vec{A}-\vec{B}}^T\brak{\vec{A}-\vec{B}} 
&+ 
\brak{\vec{B}-\vec{C}}^T\brak{\vec{B}-\vec{C}} 
\nonumber \\
&= 
\brak{\vec{C}-\vec{A}}^T \brak{\vec{C}-\vec{A}}
\nonumber \\
\implies 
2\vec{A}^T\vec{B} - 2\vec{B}^T\vec{B}&+2\vec{B}^T\vec{C}-2\vec{A}^T\vec{C}
=0
\end{align}
which can be simplified to obtain \eqref{eq:orth}.
\item Let $\vec{x}$ be any point on $AB$ in Fi.g \ref{fig:orth}.  Show that
\begin{equation}
\brak{\vec{x}-\vec{A}}^T\brak{\vec{B}-\vec{C}} = 0
\end{equation}
%
\item If $\vec{x,y}$ are any two points on $AB$, show that 
\begin{equation}
\label{eq:orth_any}
\brak{\vec{x}-\vec{y}}^T\brak{\vec{B}-\vec{C}} = 0
\end{equation}

%\numberwithin{equation}{enumi}

%\renewcommand{\theequation}{\theenumi}
\end{enumerate}
.
